\section*{Chapitre \ref{ch:m2:credit-indices-and-tranches} --- Indices de crédit et tranches}

\begin{solution}{exo:m2:credit-indices-and-tranches:1}
La part du nom vaut $1\text{ milliard}/125 = \text{USD}~8$ millions, sur laquelle on paie
$(100 - 40)\% $~: USD~4.8 millions. L'indice se poursuit sur 124 noms, pour un notionnel de
USD~992 millions, et la prime est payée sur ce montant.
\end{solution}

\begin{solution}{exo:m2:credit-indices-and-tranches:2}
0~\% (tant que les pertes restent sous 3~\%), 50~\% ($(5-3)/4$) et 100~\% (au-delà de 7~\%).
\end{solution}

\begin{solution}{exo:m2:credit-indices-and-tranches:3}
Sa liste reflète les noms investment grade liquides d'aujourd'hui, et chacun y cote et s'y
couvre, si bien que sa fourchette est la plus serrée. L'ancienne série conserve sa liste, y compris les noms
qui ont été déclassés ou sont devenus illiquides, et elle se traite de moins en moins, en off-the-run.
\end{solution}

\begin{solution}{exo:m2:credit-indices-and-tranches:4}
Non~: 146.0 points de base. Le nom le plus large a l'annuité la plus faible, 4.085 contre 4.420,
donc il pèse moins dans l'upfront moyen.
\end{solution}

\begin{solution}{exo:m2:credit-indices-and-tranches:5}
Une corrélation plus élevée fait se regrouper les défauts~: la plupart du temps, peu de défauts, voire aucun, surviennent, ce qui épargne la \omterm{def:m2:credit-indices-and-tranches:tranche}{tranche} equity, et
parfois beaucoup en surviennent, ce qui la fait perdre entièrement de toute façon. Sa perte espérée
diminue quand la corrélation augmente, si bien que le \omterm{def:m2:credit-default-swaps:cds}{vendeur de protection} y gagne.
\end{solution}

\begin{solution}{exo:m2:credit-indices-and-tranches:6}
Elle s'exerce à l'échéance si le spread de l'indice est alors supérieur à 80 points de base~: elle
achète alors une protection sur l'indice aux conditions du strike, valant environ (spread moins
80 points de base) multipliés par l'annuité de l'indice, et elle peut conserver la protection ou
la revendre au marché.
\end{solution}

\begin{solution}{exo:m2:credit-indices-and-tranches:7}
Pour $\rho = 0.3$~: 59.57~\%, 24.09~\%, 7.828~\% et 0.224~\% des \omterm{def:m2:credit-indices-and-tranches:tranche}{tranches}~; pondérées par leurs
largeurs 3, 4, 8 et 85~\%, elles totalisent 3.568~\%, soit le $5.95\% \times 60\%$ du portefeuille.
La perte espérée de la mezzanine vaut 24.3~\% pour $\rho = 0.10$, et de nouveau pour $\rho = 0.28$.
\end{solution}

\begin{solution}{exo:m2:credit-indices-and-tranches:8}
On compare l'indice à son \omterm{def:m2:credit-indices-and-tranches:intrinsic}{spread intrinsèque}, et non à la moyenne simple, qui est
plus élevée (ici de 1.7 points de base). La transaction coûte une demi-fourchette sur 125 noms et
l'indice, à l'entrée et à la sortie~; elle n'est pas sans risque, car le skew peut s'élargir avant
qu'elle ne se referme, ce qui déprécie la position et déclenche un appel de marge, et le fonds peut
ne pas pouvoir la conserver jusqu'à sa clôture. Le risque de défaut est couvert~; le skew ne l'est pas.
\end{solution}

\begin{solution}{pb:m2:credit-indices-and-tranches:1}
\textbf{1.} 75.3 et 55 points de base.
\textbf{2.} $-1.157\%$, un \omterm{def:m2:credit-indices-and-tranches:intrinsic}{spread intrinsèque} de 73.6 points de base.
\textbf{3.} On moyenne les upfronts, et non les spreads, et les noms dont le spread est large ont des
annuités plus faibles, donc ils pèsent moins.
\textbf{4.} Également 73.6~: la formule au premier ordre est très proche.
\textbf{5.} $65.6 - 73.6 = -8$ points de base.
\textbf{6.} Il reçoit 1.513~\% sur l'indice, soit USD~15.13 millions, et paie 1.157~\% sur les
composants, soit USD~11.57 millions.
\textbf{7.} USD~3.55 millions reçus~: la valeur du skew aux prix d'entrée,
que le fonds rend s'il clôture la transaction aux mêmes prix.
\textbf{8.} Il reçoit USD~4.8 millions sur l'indice et en paie USD~4.8 millions sur son
contrat sur un seul nom~: solde net nul.
\textbf{9.} USD~3.55 millions.
\textbf{10.} Une perte de USD~5.37 millions.
\textbf{11.} USD~0.657 million sur les composants et 0.109 sur l'indice, soit USD~0.77
million pour entrer~; USD~1.53 millions pour y entrer puis en sortir.
\textbf{12.} USD~2.02 millions. Le P\&L correspond à la variation du skew multipliée par l'annuité,
si bien que le skew doit varier de plus de $2 \times (1.5 + 0.25) = 3.5$ points de base.
\textbf{13.} Parce qu'un grand intervenant continue de vendre de la protection sur l'indice, les noms
individuels sont coûteux et lents à négocier, et le capital d'arbitrage est limité et doit supporter la
valorisation au prix de marché.
\textbf{14.} La marge de variation sur la position d'indice, l'usage du capital et les limites
de financement, et peut-être une sortie forcée au pire moment.
\textbf{15.} Ses ventes éloignent l'indice de ses composants, et les traders de skew qui
prennent l'autre côté en tirent profit quand il doit cesser de vendre ou racheter~; plus la
position est importante par rapport au volume du marché, plus sa sortie coûte cher.
\textbf{16.} Dénouer la position revient à racheter une protection équivalant à 10 à 15 jours du volume
total du marché, ce qui ferait fortement bouger l'indice à son encontre~; tous ceux qui suivent
le skew l'anticiperaient.
\textbf{17.} Par la perte qu'elle peut supporter si le skew double avant sa clôture, par le
volume quotidien du marché sur l'indice et les noms, et par la marge qu'elle peut financer.
\textbf{18.} La \omterm{def:m2:credit-indices-and-tranches:index}{série d'indice} passe en off-the-run~: la liquidité migre vers la nouvelle série, et
le skew de l'ancienne peut se refermer plus lentement ou être plus difficile à traiter.
\textbf{19.} Résultat nommé~: \emph{le trade sur le skew} rapporte USD~3.55 millions brut et
USD~2.02 millions nets de coûts lorsque le skew se referme de $-8$ à 0, et il perd USD~5.37
millions sur le papier s'il s'élargit d'abord jusqu'à $-20$.
\textbf{20.} Parce que la clôture dépend de l'arrêt de la contrepartie et de la capacité du fonds
à survivre à l'élargissement jusqu'à ce que la contrepartie s'arrête.
\end{solution}

\begin{solution}{iq:m2:credit-indices-and-tranches:1}
Un \omterm{def:m2:credit-default-swaps:cds}{credit default swap} (CDS) sur 125 noms nord-américains liquides et investment grade, à
pondération égale, avec un coupon de 100 points de base et coté en spread. Tous les six mois, les
20 mars et 20 septembre, une nouvelle série dont la liste est mise à jour est publiée et devient le
contrat on-the-run~; l'ancienne série se poursuit avec sa liste.
\iqlookfor{composants, coupon, le roll.}
\end{solution}

\begin{solution}{iq:m2:credit-indices-and-tranches:2}
Le spread implicite par les prix individuels des composants~: on moyenne leurs upfronts au
coupon de l'indice, puis on reconvertit. L'indice se traite selon sa propre offre et demande, est
bien plus liquide que les noms, et arbitrer entre eux est coûteux et risqué, si bien que les
deux peuvent diverger, parfois fortement.
\iqlookfor{moyenne des upfronts, et limites de l'arbitrage.}
\end{solution}

\begin{solution}{iq:m2:credit-indices-and-tranches:3}
Le nom $i$ fait défaut si $\sqrt\rho Z + \sqrt{1-\rho}\varepsilon_i < N^{-1}(p)$. Conditionnellement à $Z$, la
probabilité de défaut vaut $N\bigl((N^{-1}(p) - \sqrt\rho Z)/\sqrt{1-\rho}\bigr)$ et les
défauts sont indépendants, donc, par la loi des grands nombres, la fraction de noms en défaut
est égale à cette probabilité~; la perte vaut $1-R$ fois cette valeur.
\iqlookfor{indépendance conditionnelle et limite du grand panier.}
\end{solution}

\begin{solution}{iq:m2:credit-indices-and-tranches:4}
La \omterm{def:m2:credit-indices-and-tranches:tranche}{tranche} equity est longue en corrélation (sa perte espérée diminue quand les défauts se regroupent), les
\omterm{def:m2:credit-indices-and-tranches:tranche}{tranches} senior sont courtes. La perte espérée de la mezzanine n'est pas monotone en corrélation, de sorte que
un même prix peut correspondre à deux corrélations, ou à aucune~; les \omterm{def:m2:credit-indices-and-tranches:tranche}{tranches} de base partent toutes de zéro, et leurs
pertes espérées décroissent de façon monotone avec la corrélation, de sorte que chaque prix
donne une seule corrélation.
\iqlookfor{sens de chaque \omterm{def:m2:credit-indices-and-tranches:tranche}{tranche}, et monotonie.}
\end{solution}

\begin{solution}{iq:m2:credit-indices-and-tranches:5}
Un portefeuille conçu comme une couverture s'est transformé en un très vaste ensemble directionnel
et de valeur relative, sans limites de notionnel~; le notionnel net a triplé en un trimestre~; les
positions ont atteint un multiple important du volume quotidien du marché, si bien qu'elles ne pouvaient
pas être dénouées sans faire bouger les prix~; les valorisations n'étaient pas contrôlées indépendamment~; et les
pertes ont d'abord été sous-estimées.
\iqlookfor{taille rapportée à la liquidité, limites, contrôle des valorisations.}
\end{solution}

\begin{solution}{iq:m2:credit-indices-and-tranches:6}
S'abonner aux cotations des noms individuels et de l'indice~; maintenir à jour l'upfront et
l'annuité de chaque composant à chaque tick, avec son poids dans chaque série à laquelle il appartient~; la
valeur intrinsèque d'une série est une somme pondérée, mise à jour incrémentalement quand un
nom bouge~; on convertit en spread avec une annuité en cache et un pas de Newton~; on publie le
skew, avec l'âge et la qualité des cotations qui le sous-tendent, et signale les noms dont les données sont périmées.
\iqlookfor{sommes incrémentales, nombreuses séries par nom, données périmées.}
\end{solution}
